\chapter{Income, Wellbeing and What Economics Measures}\label{ch:b3:income-and-wellbeing-evidence}

The series has measured living standards by income: GDP per person, consumption,
wages. People can also be asked how their lives are going, and the answers can be
set against their incomes. Do richer countries report better lives? Do people
inside a country report better lives as it grows? Does the relation flatten, so
that beyond some income money stops helping? And can income be corrected for what
it leaves out --- longer lives, more leisure, a more equal spread --- in a single
number? This last chapter takes the evidence on these questions, and then looks back
at the series as a whole: which of its models explained the charts it opened with,
and which questions the evidence still leaves open.

\begin{openingchart}{Mean life satisfaction (0 to 10, people aged 16 and over) against GDP per person
(purchasing-power parity, constant 2021 dollars, logarithmic scale), 32 European countries: each
arrow runs from 2013 (circle) to 2024 (dot).}{Eurostat (ilc\_pw01), from the EU statistics on
income and living conditions; World Bank, World Development Indicators}
\begin{tikzpicture}
\begin{axis}[omscatter, width=0.72\linewidth, height=0.46\linewidth, xmin=4.2, xmax=5.16,
  ymin=4.5, ymax=8.4, xtick={4.30103,4.60206,4.90309}, xticklabels={20{,}000, 40{,}000, 80{,}000},
  xlabel={GDP per person, dollars}, ylabel={Life satisfaction, 0--10}]
  \addplot[black!45, -stealth, no marks, quiver={u=\thisrow{dx}, v=\thisrow{dy}}]
    table[x=x, y=y] {figdata/bachelor-3/36-income-and-wellbeing-evidence/out/sat_arrows.csv};
  \addplot[only marks, mark=o, mark size=1.4pt, black!55] table[x=lx13, y=s13]
    {figdata/bachelor-3/36-income-and-wellbeing-evidence/out/sat_2013.csv};
  \addplot[only marks, mark=*, mark size=1.4pt, omDef] table[x=lx24, y=s24]
    {figdata/bachelor-3/36-income-and-wellbeing-evidence/out/sat_2024.csv};
\end{axis}
\end{tikzpicture}
\end{openingchart}

% figdata: sat_stats, sat_table
Two facts sit in the chart. Across countries, the richer report more satisfied
lives, in both years. Inside countries, the arrows point up and to the right where
income grew fast from a low level --- Bulgaria from 4.8 to 6.2, Serbia from 5.0 to 6.5 --- and flat or slightly down in
the rich countries where income grew slowly: Denmark from 8.0 to 7.6, Sweden from 7.9 to 7.4. Satisfaction rose in 20 of
the 32 countries. The chapter asks what these facts mean.

\begin{recall}
Income and wellbeing are first compared in the grades 10--12 volume; what GDP leaves
out is the grades 6--9, grades 10--12 and Year 1 volumes'. The \omterm{def:b3:credible-empirical-work:evidence-scale}{evidence scale} used by
every \emph{debate} box is \cref{ch:b3:credible-empirical-work}'s; the \omterm{def:b3:cost-benefit-analysis:vsl}{value of a
statistical life} is \cref{ch:b3:cost-benefit-analysis}'s.
\end{recall}

\begin{notation}[Symbols of this chapter]\label{not:b3:income-and-wellbeing-evidence:symbols}
$S$ is mean \omterm{def:b3:income-and-wellbeing-evidence:life-satisfaction}{life satisfaction}; $y$ income per person; $c$ consumption per person; $e$ life
expectancy; $h$ hours worked per person as a share of 5{,}840 a year; $\sigma$ the standard
deviation of log consumption; $\lambda$ \omterm{def:b3:income-and-wellbeing-evidence:welfare}{consumption-equivalent welfare}.
\end{notation}

\section{Life satisfaction and income}\label{sec:b3:income-and-wellbeing-evidence:satisfaction}

\begin{definition}[Life satisfaction]\label{def:b3:income-and-wellbeing-evidence:life-satisfaction}
\index{life satisfaction}\index{Cantril ladder}
\emph{Life satisfaction} is a person's answer, on a scale usually from 0 to 10, to a
survey question asking them to evaluate their life as a whole. In the \emph{Cantril
ladder} version, 0 is the worst possible life for them and 10 the best. It is a
measure of \emph{life evaluation}, distinct from \emph{emotional wellbeing}: the
feelings of a day, asked about separately.
\end{definition}

\begin{measure}[How life satisfaction is asked]
The European survey of the opening chart asks adults to rate their satisfaction
with their life on a 0--10 scale; world surveys use the ladder. Three features shape
what the answers can show. The scale is bounded: an average of 7.6 cannot rise
much. Scale use may differ: two people with the same life may pick different numbers,
and if that differs between countries, part of a gap is a difference in language or
habit. And the question is short and depends on its place in the questionnaire. The
answers are informative nonetheless: in the same survey, people with a degree reported
7.6 on average across the European Union in 2024, people with at most lower secondary schooling 6.8.
% figdata: sat_stats
\end{measure}

\begin{definition}[Easterlin paradox]\label{def:b3:income-and-wellbeing-evidence:easterlin}
\index{Easterlin paradox}
The \emph{Easterlin paradox} is the claim that, while richer people and richer
countries report higher satisfaction at a point in time, a country's average
satisfaction does not rise as its income grows over the long run.
\end{definition}

In the opening chart's 32 countries, one more unit of log income per person is associated with 0.69 more
points of satisfaction across countries in 2024, against 1.53 in 2013. Within countries, the change in satisfaction
from 2013 to 2024 is associated with the change in log income with a slope of 1.83. One
arrow runs almost flat from about 60{,}000 to 119{,}000 dollars: Ireland, where much of GDP is
profit generated by multinational firms that goes straight to their owners abroad;
its statistics office publishes a modified national income that leaves these effects
out. % ledger: b3iw:cso.gni
Without Ireland the slope is 2.74. But
the change is strongly tied to the 2013 level (correlation $-0.87$): the countries that
started lowest rose most. Eleven years is about the shortest span the paradox speaks
of, and 2013 came at a low point for several of the countries in the chart, so a
recovery is mixed into the within-country slope.

\begin{proposition}[Equal percentage raises, equal gains]\label{prop:b3:income-and-wellbeing-evidence:log-linear}
If $S = a + b\ln y$, a raise of a given percentage $x$ raises satisfaction by $b\ln(1 + x)$ at every income; the
gain from one more dollar, $b/y$, falls as income rises. A relation that flattens when
satisfaction is plotted against income in dollars is a straight line against log
income: flattening in levels is not flattening in logs.
\end{proposition}

\begin{proof}
$S(y(1 + x)) - S(y) = b[\ln y + \ln(1 + x) - \ln y] = b\ln(1 + x)$, free of $y$. And $dS/dy = b/y$, which falls in $y$. A satiation point
would require the slope against $\ln y$ itself to fall to zero, which is a different and
testable claim.
\end{proof}

\section{Does the relation flatten?}\label{sec:b3:income-and-wellbeing-evidence:flatten}

\begin{debate}[Does wellbeing stop rising with income?]
\textbf{No long-run relation within countries.} Over ten years or more, happiness
does not rise as a country's income rises, in developed and developing countries and
in the transition economies; over shorter spans it rises in expansions and falls in
contractions (\emph{Grade 3}). % ledger: b3iw:emssz
\emph{Weakness:} short series from surveys whose questions changed, and few
countries.

\textbf{A positive relation everywhere.} With more countries and many datasets,
average satisfaction rises with GDP per person across countries with no satiation
point, the slope is similar to the one within countries, and growth over time goes
with rising satisfaction (\emph{Grade 3}). % ledger: b3iw:sw
\emph{Weakness:} cross-country correlations; the within-country series are short.

\textbf{Feelings plateau, evaluations do not.} In more than 450{,}000 responses to a daily U.S. survey of 1{,}000
residents, life evaluation on the \omterm{def:b3:income-and-wellbeing-evidence:life-satisfaction}{Cantril ladder} rises steadily with log income;
emotional wellbeing rises too, but not beyond about 75{,}000 dollars a year (\emph{Grade 3}). % ledger: b3iw:kd
\emph{Weakness:} the emotion questions were yes or no about the day before.

\textbf{No plateau.} In 1{,}725{,}994 reports sampled during the day from 33{,}391 employed U.S. adults, both
feelings and evaluations rise linearly with log income, as steeply above 75{,}000 dollars
as below (\emph{Grade 3}). % ledger: b3iw:kil
\emph{Weakness:} employed people who use a phone application, a selected group.

\textbf{A conflict resolved.} The two sets of authors reanalysed the second dataset
together. The flattening holds only for the least happy people; among happier people,
happiness rises steadily with log income, and faster among the happiest. The earlier
plateau came from yes-or-no questions that could not tell degrees of happiness apart
(\emph{Grade 3}). % ledger: b3iw:kkm
\emph{Weakness:} still a correlation of income with wellbeing, not an effect of
income on it.
\end{debate}

Across these studies, the relation between income and evaluated wellbeing is
close to log-linear and does not stop; whether a country's \emph{growth} raises its
average over decades remains disputed, because the series that could answer it are
short. Hedonic adaptation --- the measured tendency of satisfaction to return part of
the way towards its earlier level after a change in circumstances --- is one reason
offered for the paradox; it is a pattern in panel data, not a law.

\section{Welfare in one number}\label{sec:b3:income-and-wellbeing-evidence:welfare}

\begin{definition}[Consumption-equivalent welfare; equivalent income]\label{def:b3:income-and-wellbeing-evidence:welfare}
\index{consumption-equivalent welfare}\index{equivalent income}
A country's \emph{consumption-equivalent welfare} relative to a reference country is
the factor $\lambda$ by which consumption in the reference country, at every date and for
everyone, would have to be multiplied to give a person there the same expected
lifetime utility as being born in the country compared, with its consumption, leisure,
mortality and inequality. More generally, a person's \emph{equivalent income} is the
income that, under reference values of the things other than income, gives them the
same welfare as their actual situation.
\end{definition}

\begin{proposition}[The Jones--Klenow decomposition]\label{prop:b3:income-and-wellbeing-evidence:jk}
Let a year of life give flow utility $u = \bar u + \ln c - v(h)$, with $v(h) = \theta h^2/2$; let consumption across people be
lognormal with mean $c$ and log variance $\sigma^2$; and let lifetime utility be $e$ times the expected flow.
Then country $i$'s welfare relative to country $0$ is
\begin{equation*}
\ln\lambda_i = \underbrace{\frac{e_i - e_0}{e_0}\,u_i}_{\text{life expectancy}} + \underbrace{\ln\frac{c_i}{c_0}}_{\text{consumption}}
\underbrace{-\,[v(h_i) - v(h_0)]}_{\text{leisure}} \underbrace{-\,\frac{\sigma_i^2 - \sigma_0^2}{2}}_{\text{inequality}},
\end{equation*}
where $u_i = \bar u + \ln c_i - v(h_i) - \sigma_i^2/2$ is country $i$'s expected flow utility.
\end{proposition}

\begin{proof}
If $\ln c \sim N(\mu, \sigma^2)$ with mean $c$, then $c = e^{\mu + \sigma^2/2}$, so $E\ln c = \ln c - \sigma^2/2$ and expected flow utility is $u = \bar u + \ln c - v(h) - \sigma^2/2$.
Multiplying the reference consumption by $\lambda$ adds $\ln\lambda$ to its flow, so $\lambda_i$ solves $e_0(u_0 + \ln\lambda_i) = e_iu_i$, that is,
$\ln\lambda_i = (e_i/e_0)u_i - u_0 = (e_i - e_0)u_i/e_0 + (u_i - u_0)$. The last bracket is the sum of the consumption, leisure and inequality terms.
\end{proof}

The life-expectancy term depends on the level of $u$, not only on differences: what
a year of life is worth must be chosen, and with log utility that choice is the
constant $\bar u$. Here $\bar u$ is set so that a U.S. year is worth 5 in units of log consumption ---
five years of consumption at U.S. levels, of the order of the values of a statistical
life in \cref{ch:b3:cost-benefit-analysis} --- and $\theta$ so that, in the United States, the disutility of
the last hour worked equals what that hour earns, with labour taking 0.6 of income. A lognormal
spread with Gini coefficient $G$ has $\sigma = \sqrt2\,\Phi^{-1}((1 + G)/2)$.

\begin{omchart}{Consumption-equivalent welfare against GDP per person, both as a share of the
United States, 2019, 115 countries (logarithmic scales), with the line of equality.}{Penn World
Table 11.0 (GDP, consumption share, hours, employment); World Bank, World Development
Indicators (life expectancy; Gini index, latest 2010--2019); author's calculation}
\begin{tikzpicture}
\begin{axis}[omscatter, width=0.6\linewidth, height=0.46\linewidth, xmin=-1.9, xmax=0.3,
  ymin=-1.9, ymax=0.3, xtick={-1.69897,-1.30103,-1,-0.69897,-0.30103,0},
  xticklabels={2, 5, 10, 20, 50, 100}, ytick={-1.69897,-1.30103,-1,-0.69897,-0.30103,0},
  yticklabels={2, 5, 10, 20, 50, 100}, xlabel={GDP per person, \% of the U.S.},
  ylabel={Welfare, \% of the U.S.}]
  \addplot table[x=lx, y=lw] {figdata/bachelor-3/36-income-and-wellbeing-evidence/out/welfare.csv};
  \draw[black!50, dashed] (axis cs:-1.9,-1.9) -- (axis cs:0.3,0.3);
\end{axis}
\end{tikzpicture}
\end{omchart}
% figdata: welfare_stats

Across 115 countries the correlation between log welfare and log GDP per person is 0.94, but 70~\% of the
countries lie below the line: relative to the United States, their welfare is lower than
their GDP. The components show why.

\begin{center}
\begin{tabular}{lcccccc}
  & GDP & Welfare & \multicolumn{4}{c}{Terms of $\ln\lambda$} \\
  & \% of U.S. & \% of U.S. & life & consumption & leisure & inequality \\ \hline
  France & 76.1 & 103.8 & 0.246 & $-0.552$ & 0.201 & 0.143 \\
  Germany & 87.8 & 101.4 & 0.154 & $-0.388$ & 0.115 & 0.132 \\
  Japan & 63.4 & 69.6 & 0.306 & $-0.768$ & $-0.024$ & 0.123 \\
  Brazil & 24.8 & 16.0 & $-0.124$ & $-1.455$ & $-0.024$ & $-0.229$ \\
  China$^\dagger$ & 25.9 & 9.7 & $-0.029$ & $-1.872$ & $-0.484$ & 0.056 \\
  India$^\dagger$ & 10.8 & 7.9 & $-0.279$ & $-2.283$ & $-0.151$ & 0.168 \\ \hline
\end{tabular}

\smallskip{\small $^\dagger$ Gini index from a consumption survey; the others, and the
United States (41.9), from income surveys.}
\end{center}
% figdata: welfare_parts

France and Germany reach the United States in welfare: they live longer, work fewer
hours per person and are more equal, which more than offsets lower consumption. China
falls from a quarter of U.S. GDP to a tenth in welfare: its consumption share is low
and hours per person are long. Brazil's high inequality costs it a fifth of its consumption.
India's inequality term, the largest in the table, is mostly a difference of measure:
its Gini of 28.8 is of consumption, which varies less than income, and it is set
against an income Gini of 41.9 for the United States. Of the 115 countries, 57 measure
inequality by consumption. % figdata: welfare_stats

\begin{measure}[Welfare in one number]
The original measure combines consumption, leisure, mortality and inequality, first
for a few countries with detailed household data and then for many. Welfare is highly
correlated with GDP per person but often far from it: Western Europe is considerably
closer to the United States, emerging Asia has caught up less, and many developing
countries are further behind; mortality matters most. % ledger: b3iw:jk
The calculation here is a simpler one --- one year, no age structure, inequality from
the Gini index --- and it reproduces those three patterns. Its weakest numbers are the
value of a year of life, which exercise 11 varies; hours: few hours per person are
counted as leisure even when they reflect unemployment, and hours compare poorly
across countries; and inequality, whose Gini indices mix consumption surveys and
income surveys, as \cref{ch:b3:development-microeconomics} warned.
\end{measure}

\section{What the series explained, and what stays open}\label{sec:b3:income-and-wellbeing-evidence:return}

The grades 6--9 volume opened with twenty-one charts and no model: living
standards, poverty, health and population, school and work, food and energy, a
connected world, and the rest. Every later chapter opened with a chart and set a
model against it. The table gives, for each theme of those charts, a model this
volume tested and what the record showed.

\begin{center}
\small
\begin{tabular}{p{0.17\linewidth}p{0.27\linewidth}p{0.47\linewidth}}
  Theme & Model (chapter) & The record \\ \hline
  Living standards & Augmented Solow (\ref{ch:b3:growth-empirics}) & Explains most of the gaps in income; convergence only conditional, or in groups alike to start with \\
  Poverty & Poverty trap (\ref{ch:b3:development-microeconomics}) & Livestock and skills moved the poorest out of poverty, gains sustained after seven years; traps rare outside remote areas \\
  Inequality & Education and technology (\ref{ch:b3:wage-structure}) & The college premium follows relative supply over a century \\
  Health and population & Malthusian (\ref{ch:b3:demographic-transition}) & Before 1500, better technology raised density, not living standards \\
  Education & Mincer equation (\ref{ch:b3:human-capital}) & About 9~\% a year of schooling on average, stable over decades, higher in poor countries \\
  Work & Competitive market and monopsony (\ref{ch:b3:minimum-wage}) & Small employment effects of moderate minimum wages; large markdowns on some platforms \\
  Food & Random-utility demand (\ref{ch:b3:binary-outcomes}) & Forecast a new line's share within a tenth of a point \\
  Trade & Ricardo, Heckscher--Ohlin, gravity (\ref{ch:b3:ricardian-trade}, \ref{ch:b3:factor-endowments}, \ref{ch:b3:gravity}) & Ricardo fits; pure factor content fails; gravity fits best of all \\
  Prices and money & Quantity theory (\ref{ch:b3:inflation-evidence}) & Slope near 1 over decades, driven by high-inflation countries \\
  Government & Median voter, taxable income (\ref{ch:b3:public-choice-voting}, \ref{ch:b3:efficiency-cost-of-taxation}) & Close elections: winners' votes do not move toward the median \\
  Energy and environment & Kuznets curve, Hotelling (\ref{ch:b3:environmental-and-resource-economics}) & Local pollutants fall with laws; resource prices rarely follow the rule \\
  Institutions & Institutions as cause (\ref{ch:b3:institutions}) & Divided countries; the size of the effect debated \\
  Wellbeing & Log-linear satisfaction (this chapter) & Close to log-linear in income; long-run growth disputed \\ \hline
\end{tabular}
\end{center}

What stays open is what this volume's debate boxes left open, each at the grade of its
strongest evidence on the scale of \cref{ch:b3:credible-empirical-work}: whether trade raises income
(\cref{ch:b3:instrumental-variables}, grade 2); how much import competition lowered wages
(\cref{ch:b3:trade-policy}, grade 2); what caused the inflation of 2021--2023 (\cref{ch:b3:central-banking},
grade 3); whether a survey can value what is not sold (\cref{ch:b3:cost-benefit-analysis}, grade 2);
immigration and local wages (\cref{ch:b3:wage-structure}, grade 2); the employment effect of a
minimum wage (\cref{ch:b3:minimum-wage}, grade 2); geography, institutions or culture
(\cref{ch:b3:institutions}, grade 2); aid and growth (\cref{ch:b3:development-microeconomics}, grade 2); and
whether growth raises wellbeing (this chapter, grade 3). For each, the boxes name the
observation that would settle it.

\section{Exercises}\label{sec:b3:income-and-wellbeing-evidence:exercises}

\begin{exercise}[$\star$]\label{exo:b3:income-and-wellbeing-evidence:1}
From the opening chart and its data, give \omterm{def:b3:income-and-wellbeing-evidence:life-satisfaction}{life satisfaction} and GDP per person in 2013 and
2024 for Bulgaria and Denmark. Which way do their arrows point?
\end{exercise}

\begin{exercise}[$\star$]\label{exo:b3:income-and-wellbeing-evidence:2}
With $S = a + 0.69\ln y$, find the change in satisfaction from a 10~\% raise, at an income of 20{,}000 and at 80{,}000 dollars, and
the change from one more dollar at each.
\end{exercise}

\begin{exercise}[$\star$]\label{exo:b3:income-and-wellbeing-evidence:3}
A country's people live 78 years, the reference country's 81. With its flow utility $u_i = 5$, find the
life-expectancy term of $\ln\lambda$ and the welfare loss it implies, in per cent.
\end{exercise}

\begin{exercise}[$\star$]\label{exo:b3:income-and-wellbeing-evidence:4}
What data would contradict hedonic adaptation to a large rise in income?
\end{exercise}

\begin{exercise}[$\star\star$]\label{exo:b3:income-and-wellbeing-evidence:5}
Rank the chapter's 115 countries by GDP per person and by welfare. Which five rise most in the
welfare ranking and which five fall most? Why?
\end{exercise}

\begin{exercise}[$\star\star$]\label{exo:b3:income-and-wellbeing-evidence:6}
Country A consumes 40{,}000 a person and works $h = 0.15$; country B consumes 36{,}000 and works $h = 0.12$; $\theta = 40$. Find B's
equivalent consumption at A's hours.
\end{exercise}

\begin{exercise}[$\star\star$]\label{exo:b3:income-and-wellbeing-evidence:7}
With the mean fixed, the Gini index of consumption falls from 0.40 to 0.30. Under lognormality, find the change in
welfare, in per cent of consumption.
\end{exercise}

\begin{exercise}[$\star\star$]\label{exo:b3:income-and-wellbeing-evidence:8}
Three countries have mean log income 9, 10 and 11 and mean satisfaction 5.0, 6.0 and 6.7. Find the least-squares slope and the
satisfaction gain from doubling income.
\end{exercise}

\begin{exercise}[$\star\star$]\label{exo:b3:income-and-wellbeing-evidence:9}
What would show that a gap in reported satisfaction between two countries comes
from a difference in how people use the scale?
\end{exercise}

\begin{exercise}[$\star\star\star$]\label{exo:b3:income-and-wellbeing-evidence:10}
Derive the decomposition of \cref{prop:b3:income-and-wellbeing-evidence:jk}, including the inequality term, and
show that $\sigma = \sqrt2\,\Phi^{-1}((1 + G)/2)$ for a lognormal distribution.
\end{exercise}

\begin{exercise}[$\star\star\star$]\label{exo:b3:income-and-wellbeing-evidence:11}
Recompute welfare relative to the United States for France, Germany, Japan, Brazil, China
and India with the value of a U.S. year of life set at 3, 5 and 7. Which countries gain as it rises, and why?
\end{exercise}

\begin{exercise}[$\star\star\star$]\label{exo:b3:income-and-wellbeing-evidence:12}
Choose two themes of the closing table. For each, say which model explains it, what
it predicts that the record could have contradicted, and what observation would
refute it now.
\end{exercise}

\section{Problem: Fifty Years of Brenland, in One Number}\label{sec:b3:income-and-wellbeing-evidence:problem}

\begin{problem}[{Weekend problem --- Fifty years of Brenland, in one number}]\label{pb:b3:income-and-wellbeing-evidence:brenland}
Brenland's GDP per person in 1975 is the one found in \cref{ch:b3:growth-empirics}'s problem; in 2025 it comes
from \code{load("accounts")}, and the 1975 consumption share is taken equal to the 2000 one. The household
survey is \code{load("households")}. Invented: life expectancy 72.0 years in 1975 and 81.0 in 2025;
hours worked per person 900 a year in 1975 and 760 in 2025; the Gini index of consumption per person 0.27 in
1975. Kestria in 2025: 27{,}700 dalens of 2015 per person, consumption share 0.60, life expectancy 79.0, hours 820, Gini
0.33. Use the chapter's formula with Brenland in 2025 as the reference, a reference year of life worth 5, and
$\theta$ set as in the chapter. % figdata: brenland_stats

\medskip\noindent\textbf{Part I --- GDP.}
\begin{enumerate}
\item Find GDP per person in 1975 and 2025, and its growth rate a year.
\item Find consumption per person in both years.
\item Find the growth rate of consumption per person.
\item Why might the two growth rates overstate or understate the change in welfare?
\end{enumerate}

\medskip\noindent\textbf{Part II --- leisure and mortality.}
\begin{enumerate}[resume]
\item Write $\ln\lambda$ for 1975 relative to 2025.
\item Find $\theta$.
\item Find $h$ in both years.
\item Find the leisure term.
\item Find the consumption term.
\item Find the 1975 flow utility $u$ and the life-expectancy term.
\item Why does the life-expectancy term, unlike the others, depend on the value of a year?
\end{enumerate}

\medskip\noindent\textbf{Part III --- inequality.}
\begin{enumerate}[resume]
\item Find the Gini index of consumption per person in 2025 from the survey, weighting each household by
its size.
\item Find $\sigma$ in 1975 and 2025.
\item Find the inequality term, and explain its sign.
\item Find $\ln\lambda$ and $\lambda$.
\item Find the growth rate of \omterm{def:b3:income-and-wellbeing-evidence:welfare}{consumption-equivalent welfare} a year.
\item Which term adds most to welfare growth beyond consumption?
\item Find the welfare growth rate when a year is worth 3 and when it is worth 7.
\end{enumerate}

\medskip\noindent\textbf{Part IV --- Brenland and Kestria.}
\begin{enumerate}[resume]
\item Find Kestria's GDP per person as a share of Brenland's.
\item Find the four terms of Kestria's $\ln\lambda$.
\item Find Kestria's welfare as a share of Brenland's.
\item Which term accounts for most of the gap beyond GDP?
\item What would life-satisfaction surveys add to this comparison, and what could
they not settle?
\item State the result: Brenland's \omterm{def:b3:income-and-wellbeing-evidence:welfare}{consumption-equivalent welfare} growth 1975--2025, against its GDP growth.
\end{enumerate}
\end{problem}
